\section{Introduction}

One of the enduring goals in computational biology is to connect quantitative descriptions across disparate spatial and temporal scales, from molecular interactions to organism-level form. Advances in each scale in isolation have been substantial, but the construction of unified frameworks that explicitly bridge the molecular and tissue scales remains a fundamental challenge \cite{goriely2017mathematics}. In the context of growth and morphogenesis, tissue-level mechanical models of continuum media have been used to great effect \cite{moulton2013morphoelastic,moulton2020morphoelastic}, while biochemical models drawing on reaction-diffusion theory have proven remarkably flexible at reproducing diverse spatial patterns \cite{turing1952chemical,gierer1972theory,meinhardt1987model,pueyo2018reactiondiffusion}. The coupling of these scales, however, has received comparatively little attention even in systems where the connection is especially natural.

Molluscan shells provide a compelling study system for cross-scale modelling of pattern formation. The shell preserves a spatiotemporal record of its own secretion by the mantle, and the iterative geometry of this secretion has attracted quantitative attention since the foundational work of Raup \cite{raup1961geometric,raup1966geometric}. Geometric models have since been extended and reviewed extensively \cite{gerber2017geometry,urdy2011recapitulation}, but the parameters of these descriptions are typically disconnected from underlying biological processes. More recent mechanics-based models treat the mantle as a soft elastic tissue attached to the rigid shell edge, providing a physical account of ornamentation such as spines, ribs, and tubercles \cite{chirat2013mechanical,moulton2015morpho,erlich2016morphomechanics}. Biochemical descriptions of shell pigmentation, by contrast, have been developed largely independently \cite{fowler1992modeling,meinhardt1995algorithmic,boettiger2009neural}. Despite the natural suggestion that structural pattern formation may be governed by an underlying chemical pre-pattern, no prior framework has, to our knowledge, explicitly coupled these two descriptions.

In this work, we address that gap by constructing a multiscale computational framework that links a two-species Activator-Substrate reaction-diffusion model at the molecular scale to a morphoelastic rod description of mantle growth at the tissue scale. To keep the framework tractable, we treat the molecular model as a generator of spatially structured growth profiles which are then passed to the mechanical model as input, and we use a Quality-Diversity algorithm, MAP-Elites \cite{mouret2015illuminating,cully2015robots}, to efficiently explore the molecular parameter space. The implementation uses the \texttt{qdpy} Python library \cite{cazenille2018qdpy}. As a testbed, we apply the framework to the gastropod \textit{Turbo sazae}, a species exhibiting striking phenotypic plasticity in spine form \cite{fukuda2017nomenclature,tanaka2018phenotypic,hamm2008environmental}.

Our contributions are fourfold. First, we build a modular multiscale framework that explicitly maps molecular parameter sets to tissue-level spine shapes. Second, we demonstrate how pre-integration of the mechanical model, combined with a Quality-Diversity exploration, makes the framework computationally practical by replacing online simulation with lookup against a pre-computed morphospace. Third, we apply the framework to a non-model species and uncover a statistically significant morphological difference between three populations of \textit{T. sazae} collected at three Japanese localities. Fourth, we trace that phenotypic difference through the framework to parametric shifts at the biochemical scale, providing a mechanistic candidate for how environmental signals may be encoded at the molecular level.

\section{Background and Related Work}

\paragraph{Theoretical morphology of shells.} The modelling of shell form traces back to Raup's seminal morphospace parameterization of the coiled shell \cite{raup1961geometric,raup1966geometric}, and has since been extended by a wide literature on geometric and generative models \cite{gerber2017geometry,urdy2011recapitulation,meinhardt1995algorithmic}. These approaches describe form through a small set of geometric parameters but remain largely disconnected from underlying biological processes.

\paragraph{Mechanical models of shell ornamentation.} A parallel line of work treats the mantle-shell interface as a growing elastic structure subject to mechanical buckling. Chirat et al.\ \cite{chirat2013mechanical} established that spine morphologies on molluscan shells can be explained by the elastic buckling of a growing mantle that becomes longer than the current shell aperture. Extensions of this framework have accounted for rib patterns in ammonites \cite{moulton2015morpho,erlich2016morphomechanics}. A common theme across these works is the requirement for a spatial pre-pattern in either growth rate or stiffness as input; the origin of that pre-pattern is not resolved within the mechanics alone.

\paragraph{Biochemical pattern formation.} Reaction-diffusion dynamics form the basis of most biochemical models of shell pigmentation. Following Turing's chemical theory of morphogenesis \cite{turing1952chemical}, Gierer and Meinhardt proposed the foundational activator-inhibitor architecture \cite{gierer1972theory}, and Meinhardt and Klingler applied this framework to molluscan shells \cite{meinhardt1987model,fowler1992modeling,meinhardt1995algorithmic}. Boettiger et al.\ \cite{boettiger2009neural} introduced a neural-inspired formulation of shell pattern formation. A specific two-species Activator-Substrate variant from the literature \cite{fowler1992modeling} forms the molecular substrate of our own framework.

\paragraph{Quality-Diversity optimisation.} MAP-Elites \cite{mouret2015illuminating} is a black-box optimisation algorithm that maintains a structured archive of diverse high-performing solutions and has found broad application in robotics and evolutionary computation \cite{cully2015robots}. Implementations in Python are available in \texttt{qdpy} \cite{cazenille2018qdpy}. In our setting, the algorithm is used to fill a two-dimensional archive indexed by the parameters of a Gaussian growth profile with molecular parameter sets that produce that profile as output.

\paragraph{Phenotypic plasticity in molluscs.} \textit{Turbo sazae} is a marine gastropod formerly confused taxonomically with \textit{Turbo cornutus} before the nomenclatural revision of Fukuda \cite{fukuda2017nomenclature}. The species displays remarkable plasticity in spine form, with populations from exposed and sheltered shores differing visibly in horn expression. This plasticity sits within a broader literature on environmentally driven shell variation in gastropods \cite{tanaka2018phenotypic,hamm2008environmental}, but quantitative morphometric analyses have typically not been linked to underlying biochemical processes.
